Die in diesem Kapitel betrachteten Arbeiten zeigen, dass LLMs und MLLMs
zunehmend in 3D-bezogenen Aufgaben eingesetzt werden, etwa für
Szenenverständnis, 3D-Grounding, Fragebeantwortung, Interaktions- und
Aufgabenplanung, Szenengenerierung sowie artikulationsbezogene
Objektmodellierung. Ihre Ausgaben lassen sich dabei zwei Gruppen zuordnen:
natürlichsprachliche Beschreibungen, wie bei 3D-LLM und PointLLM, sowie neu
erzeugte oder veränderte Inhalte wie vollständige Szenen, Unity-Code oder
URDF-Dateien, wie bei LLMR, Holodeck, ArtLLM, ArtiWorld und URDFormer.

Trotz dieser Fortschritte bleibt ein Aspekt bislang unzureichend adressiert:
die Ableitung einer maschinenlesbaren, zustandsbasierten
Interaktionsspezifikation aus bereits vorhandenen, geometrisch unveränderten
3D-Objekten. Die vorliegende Arbeit ist somit nicht als Verfahren zur
Geometrie- oder Szenengenerierung einzuordnen, sondern als
semantisch-verhaltensbezogener Ansatz, der vorhandene Geometrie erhält und
ausschließlich um Interaktionslogik ergänzt.

Die Forschungslücke ergibt sich dabei aus der Kombination mehrerer
Anforderungen, die einzeln in verwandten Arbeiten auftreten, in dieser
Verbindung jedoch nicht adressiert werden. Vorhandene statische 3D-Objekte
sollen erhalten bleiben, relevante Objekte und Teilobjekte semantisch
interpretiert und gewünschtes Verhalten natürlichsprachlich beschrieben
werden. Zugleich soll daraus eine maschinenlesbare, Vivian-kompatible
Interaktionslogik entstehen. Da LLM-gestützte Generierungsprozesse
fehlerhafte Objektverweise, inkonsistente Strukturen oder unzulässige
Spezifikationselemente erzeugen können, sind außerdem Validierungs- und
Korrekturmaßnahmen erforderlich.

Aus dieser Forschungslücke sowie den Forschungsfragen F1 bis F4 (siehe 1.2)
folgt die Notwendigkeit eines Systems, das Szeneninformationen analysiert,
natürlichsprachliche Interaktionsbeschreibungen verarbeitet, daraus relevante
Objekte, Teilobjekte und Rollen ableitet, Vivian-kompatible Spezifikationen
erzeugt und Fehler prüft sowie gezielt korrigiert.